\chapter{Реализация модульных тестов} \label{ch4}

В данной главе описана практическая реализация модульных тестов, разработанных по принципам TDD для консольного приложения анализа палиндромных подстрок.
После доработки тестового набора проект содержит 7 test suite и 64 теста, охватывающих вычислительное ядро, эталонный алгоритм сравнения, сервисный слой, JSON-ввод, файловый вывод, assumption-style helper и консольный интерфейс.

\section{Структура тестового набора} \label{ch4:sec1}

Тестовый набор разделён на следующие группы:
\begin{itemize}
    \item \texttt{manacherOdd.test.js} --- проверка нечетных радиусов;
    \item \texttt{manacherEven.test.js} --- проверка четных радиусов;
    \item \texttt{centerExpansion.test.js} --- проверка эталонного конкурирующего алгоритма;
    \item \texttt{palindromeService.test.js} --- проверка сервисного слоя и восстановления ответа;
    \item \texttt{jsonReader.test.js} --- проверка загрузки входных данных;
    \item \texttt{resultWriter.test.js} --- проверка вывода в текстовый и JSON-формат;
    \item \texttt{consoleUi.test.js} --- проверка пользовательского консольного сценария.
\end{itemize}

Такое разбиение соответствует архитектуре программы и делает возможной изолированную проверку каждого слоя.
Дополнительно эталонный алгоритм \texttt{centerExpansion} используется при тестировании основного решения на алгоритме Манакера, что напрямую соответствует требованию задания о сравнении результатов двух разных алгоритмов.

\section{Параметризованные тесты} \label{ch4:sec2}

Для серий однотипных примеров используются параметризованные тесты Jest через \texttt{test.each()}.
Так реализованы проверки массивов радиусов для строк \texttt{"a"}, \texttt{"ab"}, \texttt{"aa"}, \texttt{"abba"}, \texttt{"abacaba"} и \texttt{"aaaaa"}.

Параметризация позволяет:
\begin{itemize}
    \item уменьшить дублирование кода;
    \item сохранить наглядность связи между входом и ожидаемым результатом;
    \item напрямую перенести в код контрольные примеры, разработанные на этапе test design.
\end{itemize}

\section{Реализация assertion и matchers} \label{ch4:sec3}

В проекте применяются как минимум пять различных способов проверки результатов:
\begin{itemize}
    \item \texttt{toEqual} --- для массивов радиусов;
    \item \texttt{toBe} --- для простых значений;
    \item \texttt{toThrow} --- для негативных сценариев;
    \item \texttt{to.equal} --- для строк и чисел;
    \item \texttt{to.deep.equal} --- для объектов результата;
    \item \texttt{to.include} --- для текстовых сообщений;
    \item \texttt{to.have.lengthOf} --- для длины массивов и списка строк.
\end{itemize}

Тем самым требование задания по assertion и по использованию разных matchers выполняется на практике.

\section{Assumption-style проверки и test doubles} \label{ch4:sec4}

В \texttt{consoleUi.test.js} добавлены assumption-style предусловия, реализованные через методы \texttt{assumeTrue} и \texttt{assumeFileExists}.
Они применяются для проверки корректности тестового окружения перед запуском сценария, связанного с help-модулем.

Для работы с внешними зависимостями используются три вида test doubles:
\begin{itemize}
    \item \texttt{spy} --- контроль вызовов \texttt{logger} и \texttt{resultPrinter};
    \item \texttt{stub} --- подмена методов файловой системы \texttt{fs.promises.access}, \texttt{fs.createReadStream} и \texttt{fs.writeFileSync};
    \item \texttt{mock} --- формальное ожидание вызова компонента записи результата.
\end{itemize}

Это позволяет отдельно проверять как сам алгоритм, так и корректность взаимодействия между слоями программы.

\section{Характерные тестовые сценарии} \label{ch4:sec5}

Ключевые примеры реализованных тестов приведены в \taref{tab:test-examples}.

\noindent
\begingroup
\centering
\small
\begin{longtable}[c]{|p{3.2cm}|p{4.7cm}|p{7.1cm}|}
    \caption{Характерные сценарии модульных тестов}
    \label{tab:test-examples}
    \\
    \hline
    Группа & Сценарий & Проверяемое требование \\ \hline
    \endfirsthead
    \captionsetup{format=tablenocaption,labelformat=continued}
    \caption[]{}\\
    \hline
    Группа & Сценарий & Проверяемое требование \\ \hline
    \endhead
    \hline
    \endfoot
    \hline
    \endlastfoot
    Нечетные радиусы & \texttt{"abacaba"} & корректность массива $d_{odd}$ \\ \hline
    Четные радиусы & \texttt{"abba"} & корректность массива $d_{even}$ \\ \hline
    Эталонный алгоритм & \texttt{expandAroundCenter("abba",1,2)} & корректность четного расширения от центра \\ \hline
    Сервисный слой & \texttt{findLongestPalindrome(...)} & восстановление наибольшего палиндрома по радиусам \\ \hline
    Рекомендации & строки разных классов & проверка всех четырёх ветвей \texttt{buildRecommendation()} \\ \hline
    JSON-ввод & файл без поля \texttt{input} & корректное исключение при неверном формате \\ \hline
    Файловый вывод & \texttt{writeResultToFile(...)} & запись JSON через stub без реального побочного эффекта \\ \hline
    Консольный UI & сценарии команд \texttt{3 -> 0}, \texttt{4 -> 0}, \texttt{5}, \texttt{7} & запуск обоих алгоритмов, повторный показ результата и вывод справки \\ \hline
\end{longtable}
\normalsize
\endgroup

\section{Результаты запуска тестов} \label{ch4:sec6}

По итоговому запуску \texttt{npm test} в проекте успешно пройдены все 64 теста.
Классическое покрытие \texttt{Jest} показало 100\% line coverage по всему проекту, 100\% statement coverage, 100\% function coverage и 96.29\% branch coverage.
Полное покрытие строк достигнуто для вычислительных модулей \texttt{analyzers}, сервисного слоя, ввода-вывода и консольного интерфейса; неполное покрытие остаётся только по отдельным ветвям значений по умолчанию в \texttt{resultWriter.js} и \texttt{consoleUi.js}.

\section{Выводы} \label{ch4:conclusion}

Реализация тестов полностью перенесла в код контрольные примеры, разработанные на этапе проектирования.
В тестовом наборе присутствуют позитивные и негативные сценарии, параметризованные проверки, assertion различных типов, два assumption-style метода, три разновидности test doubles и сравнение результата Манакера с конкурирующим алгоритмом.
Тем самым тестовая часть проекта соответствует формальным требованиям задания и при этом остаётся привязанной к реальному исполняемому коду.
